\noindent \textbf{Conventions.} All algebras $A$ in monoidal categories are assumed to be unital associative. Our~$H$ is a Hopf algebra over some fixed algebraically closed field $k$, of characteristic $0$, and $\vect$ denotes the category of $k$-vector spaces. For the purposes of this paper we are only interested in finite dimensional Hopf algebras. We~use the following version of Sweedler's notation: For~$h\in H$ we denote the coproduct $\Delta(h)\in H\ot H$ by $h^1\ot h^2$. The letter $S$ denotes the antipode of~$H$. The number $p$ is prime. Finally, DG stands for differential graded.

\section{Twisted Drinfeld double}
Let $H$ be a Hopf algebra. From~\cite{chern} we see that the study of the Hochschild and cyclic homologies of $\hmod$, the monoidal category of $H$-modules, reduces to the study of modules over a certain monad on $\hmod$. Recall that the consideration of Hochschild and cyclic homologies of monoidal categories is motivated by their recently discovered role~\cite{chern} in the understanding of Hopf-cyclic theory coefficients.

Briefly, we have the monad (see~\cite{chern} for more details):
\begin{gather}\label{themonad}
\mathop{\rm Hom}\nolimits_k(H,-)\colon \ \hmod\to\hmod
\end{gather}
with the $H$-module structure on $\mathop{\rm Hom}\nolimits_k(H,V)$:
\begin{gather}
x\cdot\varphi=x^2\varphi\big(S\big(x^3\big)(-)x^1\big),
\end{gather}
for $x\in H$ and $\varphi\in \mathop{\rm Hom}\nolimits_k(H,V)$. The unit $1_V\colon {\rm Id}(V)\to \mathop{\rm Hom}\nolimits_k(H,V)$ is
\begin{gather}
1_V(v)(h)=\epsilon(h)v
\end{gather}
and a crucial central element $\sigma_V\colon {\rm Id}(V)\to \mathop{\rm Hom}\nolimits_k(H,V)$ is
\begin{gather}
\sigma_V(v)(h)=hv.
\end{gather}

The anti-Yetter--Drinfeld contramodules then coincide with modules over this monad (shown in~\cite{chern}), while the stable ones consist of those for which the action of $\sigma$ agrees with that of $1$, and the mixed ones introduced in~\cite{chern} are the homotopic version of this on the nose requirement.

Recall the mixed complexes of~\cite{kassel}. These are complexes of vector spaces $(V^\bullet,d)$ with a~homo\-topy $h$ such that $dh+hd=0$. We~can replace vector spaces with $R$-modules for some ring~$R$. Note that as observed in~\cite{kassel}, the DG-category of mixed complexes in $R$-modules is isomorphic to the DG-category of DG-modules over $R[\theta]$, where $\theta$ is a freely and centrally adjoined degree~$-1$~graded commutative variable (naturally $R$ itself is placed in degree $0$, so that $R[\theta]=R\to R$ as a~comp\-lex) and $d=0$ on $R[\theta]$. The action of $\theta$ gives the homotopy $h$.

We can generalize the considerations of~\cite{kassel} so as to apply to our particular situation. Namely, let $z\in Z(R)$, i.e., $zr=rz$ for all $r\in R$. Define a DG-algebra $R[\theta]$ by placing $R$ in degree $0$ and $\theta$ in degree $-1$. Let $\theta$ commute with $R$ and itself, so in particular $\theta^2=0$. So far it is as above. Now define the differential to be $0$ on $R$ and $d\theta=z$. This is well defined and unique by the Leibniz rule. We~observe that the category of DG-modules over $R[\theta]$ consists of complexes of $R$-modules equipped with a homotopy $h$ such that $dh+hd=z$.

Now recall from~\cite{chern}:

\begin{Definition}
We say that $(M^\bullet,d,h)$ is a mixed anti-Yetter--Drinfeld contramodule if $(M^\bullet,d)$ is a complex of contramodules, i.e., modules over the monad \eqref{themonad}, and $h$ is a homotopy annihilating $\sigma-1$. More precisely,
\begin{gather}
dh+hd=\sigma-1.
\end{gather}
\end{Definition}

In this section we will define an explicit DG-algebra that will yield the mixed anti-Yetter--Drinfeld (aYD) contramodules (for $H$ finite dimensional) as its DG-modules. The construction of the twisted convolution algebra below is analogous to the classical Drinfeld double $D(H)$ and its anti-version $D_a(H)$~\cite{HKRS2} (we review these in Appendix~\ref{appx}, where we expand upon this comparison).

\begin{Definition}\label{dhhat}
Let $H$ be a Hopf algebra with an invertible antipode $S$, define a twisted double~$\widehat{D}(H)$ as follows. The multiplication on $\widehat{D}(H):={\rm End}(H)$ is
\begin{gather}\label{mult}
(f\star g)(h)= f\big(h^1\big)^2g\big(S\big(f\big(h^1\big)^3\big)h^2f\big(h^1\big)^1\big),
\end{gather}
thus the multiplicative identity, which we denote by $1$ is $\epsilon(-)1$, and the central element $\sigma(h)=h$ is invertible with inverse $S^{-1}$.
\end{Definition}

Definition~\ref{dhhat} is extracted from the monad~\eqref{themonad} with the sole purpose consisting of making the following lemma a tautology.

\begin{Lemma}\label{bas}
Let $H$ be a finite dimensional Hopf algebra.
\begin{itemize}\itemsep=0pt
\item The category of anti-Yetter--Drinfeld contramodules over $H$ is isomorphic to $\widehat{D}(H)$-mo\-dules.

\item The category of stable anti-Yetter--Drinfeld contramodules is isomorphic to modules over
\begin{gather}
A:=\widehat{D}(H)/(\sigma-1).
\end{gather}

\item The DG-category of mixed anti-Yetter--Drinfeld contramodules is isomorphic to DG-mo\-dules over the DG algebra
\begin{gather}
B:=\widehat{D}(H)[\theta],
\end{gather}
where $\theta$ is a freely adjoined degree $-1$ graded commutative variable and $d\theta=\sigma-1$, with $d|_{\widehat{D}(H)}=0$.
\end{itemize}
\end{Lemma}
\begin{proof}
In the finite dimensional case, as vector spaces, $\widehat{D}(H)={\rm End}(H)\simeq H^*\ot H$. Furthermore, as an algebra, $\widehat{D}(H)$ is the quotient of the free product algebra, generated by $H^*$ and $H$, by the relation:
\begin{gather}
\label{mult1}
h \chi= \chi\big(S\big(h^3\big)(-)h^1\big)h^2,
\end{gather}
where $h\in H$, $\chi\in H^*$. Thus, modules over the algebra are both $H$-modules and $H$-cont\-ra\-mo\-du\-les (same as $H^*$-modules for $H$ finite dimensional). The two actions satisfy the requisite compatibility condition for contramodules, as specified in~\cite{contra}, and ensured by \eqref{mult1}. Clearly, modules over $A$ consist of the full subcategory of objects on which $\sigma$ acts by identity, these are exactly the stable contramodules. Finally, a DG-module over $B$ is just a complex of $\widehat{D}(H)$-modules with a homotopy given by the action of $\theta$. The condition $d\theta=\sigma-1$ ensures that $dh+hd=\sigma-1$ on~$M^\bullet$.
\end{proof}

Our main goal is to compare the category of mixed
aYD contramodules to the category of~mixed complexes of stable aYD
contramodules. By the preceding lemma this means determining when,
and more interestingly when not, the category of DG-modules over $B$ is DG-equivalent to $A[\theta]$-modules (with $\theta$ of degree $-1$ and $d=0$). The study of Hopf-cyclic cohomology has thus far only concerned itself with the latter.

The following simple lemma takes care of a lot of cases.

\begin{Lemma}\label{diag}
Let $H$ be a finite dimensional Hopf algebra and suppose that the action of $\sigma-1$ on~$\widehat{D}(H)$ is diagonalizable. Then the categories of mixed complexes in stable $aYD$-contramodules and mixed $aYD$-contramodules are $DG$-equivalent.
\end{Lemma}

\begin{proof}
Since the action of the central element $\sigma-1$ on $\widehat{D}(H)$ is diagonalizable we may decom\-pose $\widehat{D}(H)$ as a product of algebras $\widehat{D}(H)_0\oplus \widehat{D}(H)_+$ with $\widehat{D}(H)_0$ being the $0$-eigenspace and~$\widehat{D}(H)_+$ all the other eigenspaces. We~have an inclusion of DGAs: $\widehat{D}(H)_0[\theta]\to B$ that induces an isomorphism on cohomology. Namely, as complexes $B$ is $\widehat{D}(H)\stackrel{\sigma-1}{\to}\widehat{D}(H)$ whereas $\widehat{D}(H)_0[\theta]$ is $\widehat{D}(H)_0\stackrel{0}{\to}\widehat{D}(H)_0$ and $\sigma-1$ is invertible on $\widehat{D}(H)_+$. Note that $\widehat{D}(H)_0\simeq A$ and we are done.
\end{proof}

